\section*{第 \ref{ch:b3:ocean-dynamics} 章——海洋动力学}

\begin{solution}{exo:b3:ocean-dynamics:1}
$f = \qty{9.37e-5}{s^{-1}}$; $\tau/(\rho f) = 0.12/(1025\times\num{9.37e-5})
= \qty{1.25}{m^2/s}$，在风的右侧，即向南。
跨越 \qty{1000}{km}，\qty{1.25e6}{m^3/s}，\qty{1.25}{\sverdrup}。
\end{solution}

\begin{solution}{exo:b3:ocean-dynamics:2}
$f = \qty{1.031e-4}{s^{-1}}$; $D_E = \pi\sqrt{2\times0.05/\num{1.031e-4}}
= \qty{98}{m}$。该处洋流为表层洋流的 $\mathrm e^{-\pi} = \num{0.043}$ 倍，
流向相反。
\end{solution}

\begin{solution}{exo:b3:ocean-dynamics:3}
$\beta = 2\Omega\cos\ang{30}/a = \qty{1.98e-11}{m^{-1}.s^{-1}}$; $V =
\num{-1e-7}/(1025\times\num{1.98e-11}) = \qty{-4.92}{m^2/s}$；跨越
\qty{5000}{km}，\qty{-24.6}{\sverdrup}：向南。
\end{solution}

\begin{solution}{exo:b3:ocean-dynamics:4}
两个\omterm{def:b3:ocean-dynamics:tides}{分潮}再次同相位所需时间为 $1/(1/12 - 1/12.4206) =
\qty{354}{h}$，\num{14.77}\text{ 天}，相当于半个朔望月。
\end{solution}

\begin{solution}{exo:b3:ocean-dynamics:5}
在北纬 \ang{40}，$f = \qty{9.37e-5}{s^{-1}}$，$\beta =
\qty{1.75e-11}{m^{-1}.s^{-1}}$，$\tau_x = \qty{0.05}{Pa}$，以及
$\partial\tau_x/\partial y = 0.1/\num{1.11e6} = \qty{9.0e-8}{Pa/m}$。
第一项为 $-\num{9.0e-8}/(1025\times\num{9.37e-5}) =
\qty{-9.38e-7}{m/s}$；第二项为 $0.05\times\num{1.75e-11}/(1025\times
(\num{9.37e-5})^2) = \qty{+9.7e-8}{m/s}$. $w_E = \qty{-8.4e-7}{m/s}$,
\qty{-26.5}{m/\text{年}}：海水被向下推。
\end{solution}

\begin{solution}{exo:b3:ocean-dynamics:6}
$\delta_S = \num{2e-6}/\num{2e-11} = \qty{100}{km}$；速度为 $=
\num{30e6}/(\num{1e5}\times800) = \qty{0.375}{m/s}$.
\end{solution}

\begin{solution}{exo:b3:ocean-dynamics:7}
$\beta = \qty{1.98e-11}{m^{-1}.s^{-1}}$; $\delta_M = (10^4/\num{1.98e-11})^{1/3}
= \qty{79.6}{km}$，急流宽度为 $3.628\,\delta_M = \qty{289}{km}$。
若 $A_H$ 减小为原来的八分之一，二者均减半，因为 $8^{1/3} = 2$。
\end{solution}

\begin{solution}{exo:b3:ocean-dynamics:8}
$f = \qty{8.57e-5}{s^{-1}}$; $v = (21.3 - 14.5)/(\num{8.57e-5}\times\num{8e4})
= \qty{0.99}{m/s}$，取正值，此时 $B$ 位于右侧：从 $A$ 望向 $B$，
朝东时，洋流流向北。海面
高出 $6.8/9.81 = \qty{0.69}{m}$（位于 $B$ 处）；
而 $B$ 处水柱较轻、较暖，位于洋流右侧。
\end{solution}

\begin{solution}{exo:b3:ocean-dynamics:9}
大潮潮差 $1.5K(1 + 0.46) = \qty{0.78}{m}$；小潮潮差 $1.5K(1 - 0.46)
= \qty{0.29}{m}$；比值 $1.46/0.54 = 2.7$。
\end{solution}

\begin{solution}{exo:b3:ocean-dynamics:10}
由 \cref{thm:b3:ocean-dynamics:kelvin} 可知，$u = \pm(g/c)\eta$ 对于 $F(x \mp
ct)$，且 $fu = -g\,\partial\eta/\partial y$ 化为 $\beta_0y(\pm g/c)G =
-gG'$，于是 $G = \exp(\mp\beta_0y^2/2c)$。仅上符号，即向东传播的波，
在赤道两侧均有界。$\beta_0 =
\qty{2.29e-11}{m^{-1}.s^{-1}}$，$\sqrt{2.5/\num{2.29e-11}} = \qty{330}{km}$。
\end{solution}

\begin{solution}{exo:b3:ocean-dynamics:11}
(a)~将 $\beta v = f\,\partial w/\partial z$ 从底部（$w = 0$）
积分至深层顶部（$w = W_0$），得 $\beta V = fW_0$，且 $f/\beta =
a\tan\varphi$。(b)~$V = \num{6.371e6}\times0.577\times10^{-7} =
\qty{0.368}{m^2/s}$，\qty{1.84}{\sverdrup}，跨越 \qty{5000}{km}，向极地。
(c)~来自由高纬度源区补给的深层西边界流；
内部流动指向源区，边界流则背离
源区。
\end{solution}

\begin{solution}{exo:b3:ocean-dynamics:12}
(a)~取 $z$ 自海底竖直向上，$A_zW'' = \mathrm if(W - W_g)$，其中
压强梯度取上方地转流的值。$W - W_g =
C\mathrm e^{-(1+\mathrm i)z/\delta}$ 向上有界，且 $W(0) = 0$ 给出
$C = -W_g$。(b)~输运亏损为 $\int_0^\infty(W - W_g)\,\mathrm dz
= -W_g\delta/(1+\mathrm i) = -\tfrac12\delta(1 - \mathrm i)W_g$，即
$-\tfrac12\delta(u_g + v_g,\ v_g - u_g)$。地转流
无辐散，$w = -\nabla\cdot(\text{亏损}) = \tfrac12\delta(\partial_x
v_g - \partial_y u_g) = \tfrac12\delta\zeta_g$。(c)~$\delta =
\sqrt{0.02/10^{-4}} = \qty{14.1}{m}$; $r = 10^{-4}\times14.1/8000 =
\qty{1.77e-7}{s^{-1}}$，旋转减速时间为 \num{65}\text{ 天}。
\end{solution}

\begin{solution}{pb:b3:ocean-dynamics:1}
\textbf{1.} $\tau_x = \qty{-0.1}{Pa}$（东风），0 和 \qty{+0.1}{Pa}
（西风）。

\textbf{2.} $f = \num{4.99e-5}$，\num{8.37e-5} 和
\qty{1.117e-4}{s^{-1}}。

\textbf{3.} $V_E = -\tau_x/(\rho f)$：在北纬
\ang{20} 处为 \qty{+1.96}{m^2/s}（向北），而北纬 \ang{50} 处为 \qty{-0.87}{m^2/s}（向南）。二者汇聚。

\textbf{4.} $D_E = \pi\sqrt{2\times0.05/\num{8.37e-5}} = \qty{109}{m}$。

\textbf{5.} $\tau_x = 0$ 处位于北纬 \ang{35}，故 $w_E = -(1/\rho f)\,
\partial\tau_x/\partial y = -\pi\tau_0/(\rho fL) = \qty{-1.10e-6}{m/s}$,
\qty{-34.7}{m/\text{年}}，向下。

\textbf{6.} $500/34.7 = \num{14.4}\text{ 年}$。

\textbf{7.} $L = a\times\pi/6 = \qty{3336}{km}$; $\beta =
\qty{1.875e-11}{m^{-1}.s^{-1}}$.

\textbf{8.} $\operatorname{curl}_z\vec\tau = -\pi\tau_0/L =
\qty{-9.42e-8}{Pa/m}$。

\textbf{9.} $V = \num{-9.42e-8}/(1025\times\num{1.875e-11}) =
\qty{-4.90}{m^2/s}$，向南。

\textbf{10.} $4.90\times\num{6e6} = \qty{2.94e7}{m^3/s}$,
\qty{29.4}{\sverdrup} 向南。

\textbf{11.} $\qty{4.9}{mm/s}$；\qty{1000}{km} 用时
\qty{2.0e8}{s}，\num{6.5}~年。

\textbf{12.} $fw_E/\beta = \num{8.37e-5}\times(\num{-1.10e-6})/
\num{1.875e-11} = \qty{-4.90}{m^2/s}$：与上式相同，因为埃克曼
输送在 $\tau_x = 0$ 处为零。以速率 $w_E$ 被压缩的水柱
其\omterm{def:b3:atmospheric-dynamics:pv}{位涡} $f/h$ 只能通过向南移动到较小的 $f$ 处来保持，
以速率 $\beta V_g = fw_E$ 进行。

\textbf{13.} 边界层只能是这样一层，其中对内部的修正
沿远离边界的方向衰减；当 $\beta > 0$ 时，该修正为
$\mathrm e^{-x/\delta_S}$，向东衰减，仅从西边界算起。
从涡度的角度看，北向水流紧贴东边界时所受的摩擦
会消去正涡度，符号不对。

\textbf{14.} $\delta_S = \num{e-6}/\num{1.875e-11} = \qty{53.3}{km}$。

\textbf{15.} $A = \qty{4.90}{m^2/s}$，$X/\delta_S = 112.5$：
$4.90\times111.5 = \qty{546}{m^2/s}$，\qty{0.55}{m/s} 穿过
\qty{1000}{m}。

\textbf{16.} 当 $(X/\delta_S)\mathrm e^{-x/\delta_S} = 1$ 时，$x =
\delta_S\ln(X/\delta_S) = \qty{252}{km}$。

\textbf{17.} $A(X - \delta_S - x_0) = 4.90\times(6000 - 53 - 252)\times
10^3 = \qty{27.9}{\sverdrup}$。紧贴边界的条带中，内部流动仍向南，
缺额 $A(x_0 + \delta_S) =
\qty{1.5}{\sverdrup}$，是本应穿过
该条带的内部输送量。

\textbf{18.} $\delta_M = (\num{5e3}/\num{1.875e-11})^{1/3} = \qty{64.4}{km}$；
急流 $3.628\times64.4 = \qty{233}{km}$。

\textbf{19.} $1.163\times29.4 - A\times\qty{233}{km} = 34.2 - 1.1 =
\qty{33.0}{\sverdrup}$；超出部分，即内部输送的 \qty{16}{\percent}，
在急流正东侧的逆流中向南回流。

\textbf{20.} $c = \sqrt{9.81\times4000} = \qty{198}{m/s}$；$L_R =
198/\num{8.37e-5} = \qty{2370}{km}$。

\textbf{21.} $198\times12.4206\times3600 = \qty{8860}{km}$。

\textbf{22.} 周长 $2(6000 + 3336) = \qty{18670}{km}$，\qty{26.2}{h}，
2.1 个 M$_2$ 周期。

\textbf{23.} 该圆周长为 $2\pi a\cos\ang{35} = \qty{32790}{km}$，
在一个太阴日（\qty{24.84}{h}）内走完，速度为 \qty{367}{m/s}，需要 $c^2/g =
\qty{13.7}{km}$ 厚的水层。

\textbf{24.} 罗斯贝半径与海盆宽度相当，
且 M$_2$ 波长与其长度相当，约有一个\omterm{def:b3:ocean-dynamics:amphidrome}{无
潮点}，位于海盆中央附近；高潮绕其逆时针
旋转，使海岸位于其右侧。

\textbf{25.} $\tfrac34\times0.357\times\cos^2\ang{35} = \qty{0.18}{m}$。

\textbf{26.} 西边界流向北回流的
\textbf{\qty[text-series-to-math=true]{29}{\sverdrup}}，即内部向南输送的量
（由计算得 \qty{29.4}{\sverdrup}）。
\end{solution}
